%  LaTeX-mal for labrapporter i fysikk, v.01.11.2011 
% F{\o}rst m{\aa} vi definere noen ting for dokumentet:
\documentclass[1p]{article}		% 5p gir 2 kolonner pr side. 1p gir 1 kolonne pr side. % fikk til aa fjerne preprint, gjorde om \documentclass[1p]{elsarticle} til \documentclass[1p]{article} og kommenterte ut journal		
%\journal{-}
\usepackage[usenames,dvipsnames]{color}
\usepackage{listings}
\usepackage{placeins}
\usepackage{booktabs}
\usepackage{float}
\lstset{ %
  backgroundcolor=\color{yellow2}
  \lstset{escapechar=\@}
  }
\usepackage{color}
\definecolor{yellow2}{RGB}{255,255,255} %{255,253,208}
\usepackage{ dsfont }
\usepackage{dcolumn}
\usepackage{enumitem} %abc
\usepackage[utf8]{inputenc}
\usepackage[T1]{fontenc} 						% Vise norske tegn.
\usepackage{graphicx}
\usepackage{caption}
\usepackage{subcaption}
\usepackage[T1]{fontenc} 
\usepackage{fancyvrb} 
\usepackage{listings} 
\usepackage{bm} 
\usepackage{xcolor} 
% \usepackage[latin1]{inputenc}		% Velger tengsettet i dette dokumentet			
\usepackage[english]{babel}								% Tilpasning til norsk.
\usepackage{graphicx}       						% For {\aa} inkludere figurer.
\usepackage{amsmath,amssymb} 						% Ekstra matematikkfunksjoner.
\usepackage{siunitx}							      % M{\aa} inkluderes for blant annet {\aa} f{\aa} tilgang til kommandoen \SI (korrekte m{\aa}ltall med enheter)
\usepackage{hyperref}
	\sisetup{exponent-product = \cdot}      % Prikk som multiplikasjonstegn (i steden for kryss).
 	\sisetup{output-decimal-marker  =  {,}} % Komma som desimalskilletegn (i steden for punktum).
\usepackage{booktabs}                     % For {\aa} f{\aa} tilgang til finere linjer (til bruk i tabeller og slikt).
%\usepackage{listings,relsize}
\lstloadlanguages{R}
\newcommand{\lil}[1]{\lstinline|#1|}
%\lstset{ backgroundcolor=\color{yellow2}
%\lstset{escapechar=\@}}  
\xdefinecolor{gray}{rgb}{0.4,0.4,0.4} 
\xdefinecolor{blue}{RGB}{58,95,205}% R's royalblue3; #3A5FCD 
%\lstset{% setup listings 
 %       language=R,% set programming language 
  %      basicstyle=\ttfamily\small,% basic font style 
   %     keywordstyle=\color{blue},% keyword style 
    %    commentstyle=\color{gray},% comment style 
     %   numbers=left,% display line numbers on the left side 
      %  numberstyle=\scriptsize,% use small line numbers 
       % numbersep=10pt,% space between line numbers and code 
        %tabsize=3,% sizes of tabs 
       % showstringspaces=false,% do not replace spaces in strings by a certain character 
        %captionpos=b,% positioning of the caption below 
       % breaklines=true,% automatic line breaking 
        %escapeinside={(}{)},% escaping to LaTeX 
     %   fancyvrb=true,% verbatim code is typset by listings 
      %  extendedchars=false,% prohibit extended chars (chars of codes 128--255) 
       % literate={"}{{\texttt{"}}}1{<-}{{$\bm\leftarrow$}}1{<<-}{{$\bm\twoheadleftarrow$}}1 
       % {~}{{$\bm\sim$}}1{<=}{{$\bm\le$}}1{>=}{{$\bm\ge$}}1{!=}{{$\bm\neq$}}1{^}{{$^{\bm\wedge}$}}1,% item to replace, text, length of chars 
        %alsoletter={.<-},% becomes a letter 
        %alsoother={$},% becomes other 
        %otherkeywords={!=, ~, $, \&, \%/\%, \%\%, \%\%, <-, <<-, /},% other keywords 
        %deletekeywords={c}% remove keywords 
%} 
%\usepackage[font=small,labelfont=bf]{caption}	% For justering av figurtekst og tabelltekst.
%%%%%%%%%%%%%%%%%%%%%%%%%%%%%%%%%%%%%%%%%%%%%%%%%%%%%%%%%%%%%%%%%%%%%%%%%
\usepackage[a4paper]{geometry}
\title{}
\author{}
\begin{document}\begin{center}\textbf{
{\huge{TMA4265 Stochastic Modelling -- Fall 2026}}\\ \vspace{3mm}
{\LARGE{  Project 2}}}\end{center}
\vspace{5mm}
{\LARGE \textbf{Background information}}\\
\begin{itemize}
  \item
    The deadline for the project is Sunday the 8th of November at 23:59.
  \item
    A compressed/zipped folder containing the PDF-file of the report and the files with computer code must be submitted via email. Send to \texttt{elling.svee@ntnu.no}, with subject line "TMA4265 Project 2". Imporantly, include the names of all group members in the email. The email must be sent BEFORE the deadline.
    \item
    This project must be passed to be admitted to the final exam. 
  \item Problem 1 and 2 are compulsory and \textbf{a reasonable attempt must be made on each of the subproblems in order to pass}.
    \item In addition, you need a total of \textbf{50 points} to pass. Problem 1 is worth up to 70 points, Problem 2 is worth up to 30 points, and Problem 3 is worth up to 20 points.
    \item Attempting Problem 3 is optional.
  \item
    The project should be done in groups of 3 people. Exceptions to this fixed group size can be made.

  \item
    The project report should preferably be prepared using Latex or Typst, and \textbf{must}:
    \begin{itemize} 
      \item
        Be a PDF-file.
      \item
        Include the equations and explanations necessary to justify the answers in each problem.
      \item 
        Include the required plots in the PDF and reference them in the text.
    \end{itemize}
    The project report should be a complete text, meaning it should be readable to someone with no knowledge of the questions. 
  \item
    The computer code must be written in \texttt{R}, \texttt{Python} or \texttt{Julia}. Using a \texttt{Jupyter Notebook} or \texttt{R Markdown} is allowed. The code must be submitted as a separate file \textbf{not} as an appendix in the report. Make sure your code runs. We may test it. Also make your computer code readable and add comments that describe what your code is doing.
  \item
    There is a \textbf{8 page limit} for the project report, \textbf{or 10 pages} if you attempt the third optional problem. If you submit a longer report,
    we may not read it. The computer code is submitted as a separate file and may be as long as necessary.
  \item
    See the wiki for information about digital sessions where you can ask questions about the project. You can also ask questions on the forum (\url{https://mattelab2026h.math.ntnu.no/c/tma4265/34}), or email the teaching assistant.
  \item
    Write a small paragraph at the end of your report describing how you divided the workload between the group members.
\end{itemize}
\newpage
\vspace{1cm}
\noindent{\LARGE{\textbf{Problem 1: Urgent care center}}}\\\\
\textbf{Each subproblem is worth up to 10 points.} \\\\
A small urgent care center (UCC) only has capacity to treat one patient at the time.
Assume that the arrival of patients to the UCC follows a Poisson process with
rate $\lambda > 0$, that the treatment times of the patients follow independent 
exponential distributions with expected value $1/\mu > 0$, and that treatment times
are independent of the arrival processes to the UCC. If no one is receiving treatment, 
an arriving patient
will immediately start receiving treatment. If another patient is
currently receiving treatment, the arriving patient will join (at the end of) the queue of 
patients waiting for treatment. For the purpose of this project, we will assume that
there is no upper limit on the number of patients that can wait inside the UCC.

Assume that $\lambda < \mu$, and let $X(t)$ denote the number of patients in the UCC at time
$t$. We know that $X(0) = 0$.\\\\
\textbf{a)} Complete the following tasks:
\begin{itemize}
  \item State the conditions for the UCC to be an M/M/1 queue, and explain why each of
  the conditions are satisfied.
  \item Explain why the stochastic process
$\{X(t): t\geq 0\}$ can be viewed as a birth-and-death process, and determine the birth rates
and the death rates.
  \item Determine the average time a patient will spend in the UCC as a function
  of $\lambda$ and $\mu$.
\end{itemize}
\textbf{b)} Assume for this subproblem that $\lambda = 5$ patients per hour and
that $1/\mu = 10$ minutes, and complete the following tasks:
\begin{itemize}
  \item (Computer code) Simulate $\{X(t): t\geq 0\}$ for 50 days.
  \item (Computer code) Based on the simulation, estimate the expected time a patient spends in the UCC. \emph{Hint: Estimate the long-run mean number of patients in the UCC, and use Little's law to find an estimate for the expected time.}
  \item (Computer code) Repeat 30 times and compute an approximate 95\% confidence interval (CI)
  for the expected time a patient spends in the UCC.
  \item (Report) Plot one realization of $\{X(t): t \geq 0\}$ for the time 0--12 hours. This plot must show a series of horizontal lines (and not one continuous graph).
  \item (Report) Briefly describe the steps you used to estimated the expected time, give the CI, and compare the CI to the exact value computed using the expression from 1a).
\end{itemize}

\noindent
The above model is unrealistic as the UCC will operate with a triage system where
patients are prioritized according to urgency upon arrival. 
Assume that at arrival each patient will be immediately classified as urgent or normal.
The probability that the patient is urgent is $0 < p < 1$, and the probability
that the patient is normal is $1-p$. We assume that the classification is independent
between patients.

The UCC operates with the following queue discipline. Normal patients will not receive treatment
until there are no urgent patients in the UCC. If a normal patient is receiving
treatment when an urgent patient arrives, the urgent patient will immediately start receiving
treatment and the normal patient will be moved back into the front of the queue. 

Let $U(t)$ and
$N(t)$ denote the number of urgent and normal patients, respectively, at time $t$. Since $X(0) = 0$, we also have that $U(0) = N(0) = 0$.
\newpage
\noindent\textbf{c)} Complete the following tasks:
\begin{itemize}
  \item Show that $\{U(t): t \geq 0\}$ satifies the conditions of an M/M/1 queue.
  \item Determine the arrival rate.
  \item Determine the long-run mean number of urgent patients in the UCC
        as a function of $\lambda$, $\mu$ and $p$.
\end{itemize}
\textbf{d)} Complete the following tasks:
\begin{itemize}
  \item
    Explain in words why $\{N(t): t \geq 0\}$ does not behave as an M/M/1 queue.\\
    \emph{Hint: In one case the number of normal patients transitions from 2 to 1, and in another, it transitions from 0 to 1. How might these two cases be different?}
  \item
    Determine the long-run mean number of normal patients in the UCC
    as a function of $\lambda$, $\mu$ and $p$. 
\end{itemize}
\textbf{e)} Use Little's law to show that:
\begin{itemize}
  \item The expected time in the UCC for an urgent patient as a function of $\lambda$, $\mu$ and
  $p$ is given by
  \[
    W_\mathrm{U} = \frac{1}{\mu-p\lambda}.
  \]
  \item The expected time in the UCC for a normal patient as a function of $\lambda$, $\mu$ and
  $p$ is given by
  \[
    W_\mathrm{N} = \frac{\mu}{(\mu-\lambda)(\mu-p\lambda)}.
  \]
\end{itemize}
\textbf{f)}
  Assume that $\lambda = 5$ patients per hour and that $1/\mu = 10$ minutes. We want to study
  how the expected time in the UCC varies as a function of the proportion of urgent
  patients that arrive, $p$. Complete the following tasks:
  \begin{itemize}
    \item
      Plot $W_\mathrm{U}$ and $W_\mathrm{N}$ as functions of $p$ in the same figure.
    \item
      Give interpretations of what the situations $p \approx 0$ and $p \approx 1$ describe.
    \item
      Calculate the expected time spent at the UCC for a normal patient, $W_\mathrm{N}$, in the extreme cases $p \approx 0$ and $p \approx 1$.
    \item 
      Calculate (by hand) the $p$ for which the expected time spent at the UCC for a normal
      patient is 2 hours.
  \end{itemize}
\textbf{g)} Assume that $\lambda = 5$ patients per hour, that $1/\mu = 10$ minutes, and that
$p = 0.80$. Complete the following tasks:
\begin{itemize}
  \item
    (Computer code) Simulate $\{U(t): t\geq 0\}$ and $\{N(t): t \geq 0\}$ jointly for 50 days.
  \item
    (Computer code) Estimate the expected time in the UCC for an urgent patient, and estimate the expected time in the UCC for a normal patient. Repeat 30 times to compute an
    approximate 95\% CI.
  \item
    (Report) Plot one joint realization of the total number of (normal and urgent) patients in the UCC and the
    number of urgent patients in the UCC for time 0--12 hours in the same figure. Colour them in different
    colours. The realizations must be displayed as series of horizontal lines (and not continuous graphs).
  \item
    (Report) Briefly describe how you estimated the two quantities, give the two CIs, 
    and compare them to the exact numbers calculated using the expressions in 1e).
\end{itemize}




\newpage 
\noindent
{\LARGE \textbf{Problem 2: Calibrating climate models}}\\\\
\textbf{Each subproblem is worth up to 10 points.} \\\\
A group of climate scientists are running a climate model that outputs
the temperature at every location on earth for every 6-hour period in
the years 2006 and 2100\footnote{See, for example, \url{http://www.cesm.ucar.edu/projects/community-projects/LENS/}}. The climate model is deterministic, and given
the atmospheric starting conditions, external forcing, and model parameters, you will always get the
same result. The challenge is that the parameters of the climate model
must be selected so that the output provides as realistic evolution in time as possible. This
is immensely difficult because running the model only once may require one month
of computation time. For the sake of this project, assume that the only way to choose these
parameters is to run the climate model for different parameter values and 
compare to observed temperatures.\\\\
We limit the focus to one parameter, ``the albedo of sea ice'', which is a measure how much
sun light is reflected by sea ice. We call this parameter $\theta$, and we decide to choose
this parameter so that the temperatures observed from January 1, 2006, to October 24, 2023,
matches the output of the climate model as well as possible. The fit is measured 
through a score $y(\theta)$ calculated based on the model output generated with parameter
value $\theta$.\\\\
The group of climate scientists have spent the last month running the model in five computing
centres and provides you with five evaluation points of $(\theta, y(\theta))$: $(0.30, 0.5)$, 
$(0.35, 0.32)$, $(0.39, 0.40)$, $(0.41, 0.35)$, and $(0.45, 0.60)$. The
observations are shown in Figure \ref{fig:ComputerExp}.\\\\ 
You will use a Gaussian process model $\{Y(\theta): \theta\in[0,1]\}$ to model the unknown
relationship between the parameter value and the score. Use 
$\text{E}[Y(\theta)] \equiv 0.5$,
$\text{Var}[Y(\theta)] \equiv 0.5^2$, and
$\text{Corr}[Y(\theta_1), Y(\theta_2)] = (1+15\vert \theta_1-\theta_2\vert)\exp(-15\vert\theta_1-\theta_2\vert)$ for $\theta_1, \theta_2 \in [0,1]$.\\\\
\begin{figure}
  \centering
  \includegraphics[width = 10cm]{observations.eps}
  \caption{Observed relationship between fit and model parameter.}
  \label{fig:ComputerExp}
\end{figure}

\noindent
\textbf{a)} Define a regular grid of parameter values from $\theta = 0.25$ to 
$\theta = 0.50$ with spacing $0.005$ ($n = 51$ grid points). Construct the mean
vector and the covariance matrices required to compute the conditional mean
and covariance matrix of the process at the 51 grid points conditional on the five
evaluation points. Display the prediction as a function of $\theta$, 
along with 90\% prediction intervals.\\[0.2cm] 
\emph{Hint: The quantiles of a Gaussian distribution can be calculated in R (\texttt{qnorm}) or Python \linebreak(\texttt{scipy.stats.norm.ppf})} \\[0.3cm]
\textbf{b)} The scientists' goal is to achieve $y(\theta) < 0.30$. Use the predictions
from \textbf{a)} to compute the conditional probability that $Y(\theta) < 0.30$ given
the 5 evaluation points. Plot the probability as a function of $\theta$.\\\\
\textbf{c)} The scientists decide to run the climate model again with 
$\theta = 0.33$ and the result is $y(\theta) = 0.40$. Add this to the set of observed
values, and given the six evaluation points, compute and visualize the predictions, 90\%
prediction intervals, and the probabilities that $Y(\theta) < 0.30$. The scientists'
budget allow for one more run of the climate model, which value of $\theta$ would you
suggest them to use to have the best chance to achieve $y(\theta) < 0.30$ and why?

\newpage 
\noindent
{\LARGE \textbf{Problem 3: Geometric Brownian motion (OPTIONAL)}}\\\\
It is optional to attempt this problem, but \textbf{each subproblem is worth up to 5 points}. You may find Section 3.2 from the Gaussian process notes helpful when solving this problem, particularly in part of part c). \\\\
Let
\begin{equation}
	X(t) = \sigma B(t), \quad t \geq 0,
\end{equation}
where B(t) is standard Brownian motion with unit variance, $\sigma > 0$ is the scale parameter, and with $X(0) = B(0) = 0$. Let $\tau_{a,b} = \min\{t\geq 0: X(t) \in \{-b, a\}\}$ for $a, b, > 0$ be the first time the stock reaches threshold $a$ or $-b$. We are interested in the probability of the process reaching threshold $-b$ given it starts at 0. The probability of reaching threshold $a$ before $b$ is given as: 
\begin{equation}
P(X(\tau_{a,b})=a \mid X(0)=0)= \frac{b}{a+b}.
\end{equation}
The expected amount of time to first reach $a$ or $b$ is:
\begin{equation}
E[\tau_{a,b}] = \frac{ab}{\sigma^2}.
\end{equation}

\noindent
Per Ivar decided to buy some cryptocurrency for a price of 50kr and decides to sell them when they reach 75kr or drop down to 25kr. Assume the currency price, $Y(t)$, follows geometric Brownian motion, with
$$ Y(t) = \exp\{\mu + \sigma B(t)\} \quad t \geq 0, $$
for $B(0) = 0,$ so that $Y(0) = e^\mu$ and $\mu = \log(50)$, and where $t$ is in days traded after Per Ivar bought the currency. Assume the scale parameter is given by $\sigma^2=4$.\\

\noindent
\textbf{a)} What is the probability that Per Ivar profits from selling the currency?\\

\noindent
\textbf{b)} Calculate by hand the expected time in days until Per Ivar will have to sell his cryptocurrency regardless of whether he will sell it at loss or profit from it. Write code in R to confirm the calculated values through simulation. Plot 10 realizations (simultaneously) of the process from time $t=0$ until $Y(t)$ is either 25 or 75 kroner.\\

\noindent
\textbf{c)} Numerically compute the expected time in days that Per Ivar owned the cryptocurrency before selling, conditional on the price reaching 75 kroner before 25 kroner, in the following ways:
\begin{itemize}
\item Using simulation in R or Python, by simulating at least 1,000 realizations of the process.
\item Using numerical integration in R or Python.\\
\emph{Hint: if $T>0$ is a positive random variable, and $S(t) = P(T \geq t)$, then $E[T] = \int_0^\infty S(t) \, \mathrm{d}t$. In \texttt{R}, you may use the \texttt{integrate} function. In \texttt{Python}, you can use \texttt{scipy.integrate}.}
\end{itemize}

\end{document}

